% ─────────────────────────────────────────────────────────────────────────────
% Supplementary: details of Sec. 3.2. tab_recall_full.tex and the numbers here come from paper/figures/make_recall_figure.py
% and results/diagnostics/vstar_selection/ (post-hoc analysis of frozen v1.0, V* Bench only).
% ─────────────────────────────────────────────────────────────────────────────
\section{Details of the reliability analysis}
\label{sec:reliability_supp}

\begin{table}[t]
  \centering
  \footnotesize
  \setlength{\tabcolsep}{4pt}
  \begin{tabular}{lc cc cc}
    \toprule
    $B$ & $K$ & Random & Uniform & Coarse prior & After scoring \\
    \midrule
    1024 & 2 & 7.2 & 3.7 & 27.7 & \textbf{69.1} \\
    1536 & 4 & 14.8 & 8.9 & 67.0 & \textbf{71.2} \\
    2048 & 4 & 14.8 & 8.9 & 59.7 & \textbf{72.8} \\
    3072 & 6 & 22.6 & 14.1 & 77.5 & \textbf{80.1} \\
    \bottomrule
  \end{tabular}
  \caption{\textbf{Recall@$K$ (\%) on V$^*$~Bench} at about the number of detail tiles each budget keeps, with $K$ fixed so that all columns are comparable. Random and uniform selection do not depend on the budget.}
  \label{tab:recall}
\end{table}

\begin{table}[t]
  \centering
  \footnotesize
  \setlength{\tabcolsep}{2.5pt}
  \begin{tabular}{lc ccc cc}
    \toprule
    $B$ & Window & Prior & Recov. & Lost & Final & Rate (rand.) \\
    \midrule
    1024 & 2 of 5 & 53 & +79 & 0 & \textbf{69.1\%} & 89.8\% (33.8\%) \\
    1536 & 2 of 4 & 126 & +17 & $-$7 & \textbf{71.2\%} & 77.3\% (39.4\%) \\
    2048 & 2 of 4 & 115 & +27 & $-$2 & \textbf{73.3\%} & 79.4\% (33.8\%) \\
    3072 & 1 of 3 & 149 & +9 & $-$4 & \textbf{80.6\%} & 64.3\% (31.0\%) \\
    \bottomrule
  \end{tabular}
  \caption{\textbf{From the prior to the final selection} on the 191 questions of V$^*$~Bench. Prior: questions on which the prior alone covers the target with the tiles of the budget. Recov.: the target enters the window, the prior alone would miss it, and boundary scoring keeps it. Lost: the reverse. Final = (prior + recovered $-$ lost) / 191. Rate: the recovery rate, computed over the questions whose target enters the window and that the prior alone would miss (88, 22, 34, 14 questions); in parentheses, the expected rate when the window tiles are kept at random.}
  \label{tab:window}
\end{table}

\begin{table*}[t]
  \centering
  \footnotesize
  \setlength{\tabcolsep}{4pt}
  \begin{tabular}{r cc cc cc cc}
    \toprule
    & \multicolumn{2}{c}{$B=1024$} & \multicolumn{2}{c}{$B=1536$} & \multicolumn{2}{c}{$B=2048$} & \multicolumn{2}{c}{$B=3072$} \\
    \cmidrule(lr){2-3}\cmidrule(lr){4-5}\cmidrule(lr){6-7}\cmidrule(lr){8-9}
    $K$ & Prior & After scoring & Prior & After scoring & Prior & After scoring & Prior & After scoring \\
    \midrule
    1 & 3.1 & 51.3 & 3.1 & 3.1 & 3.1 & 3.1 & 3.1 & 3.1 \\
    2 & 27.7 & 69.1 & 26.7 & 30.4 & 23.6 & 23.6 & 27.7 & 27.7 \\
    3 & 51.3 & 71.2 & 50.3 & 61.3 & 46.1 & 61.8 & 42.9 & 42.9 \\
    4 & 68.6 & 72.8 & 67.0 & 71.2 & 59.7 & 72.8 & 58.1 & 58.1 \\
    5 & 73.8 & 73.8 & 72.8 & 74.3 & 69.1 & 77.0 & 68.1 & 68.1 \\
    6 & 78.5 & 78.5 & 78.5 & 78.5 & 77.5 & 78.0 & 77.5 & 80.1 \\
    8 & 85.3 & 85.3 & 85.9 & 85.9 & 84.3 & 84.3 & 85.3 & 85.3 \\
    12 & 95.3 & 95.3 & 94.8 & 94.8 & 92.1 & 92.1 & 93.2 & 93.2 \\
    \bottomrule
  \end{tabular}
  \caption{\textbf{Recall@$K$ (\%) of the coarse prior and after boundary scoring} on V$^*$~Bench, at every budget with saved selections.}
  \label{tab:recall_full}
\end{table*}


\cref{tab:recall,tab:window} give the numbers behind \cref{sec:reliability}. Little has to be found: of a median of
24 detail tiles, one tile covers the target on 78.5\% of the questions, and two tiles on all of them. Boundary scoring
recovers the targets that the prior alone would miss at 2.0 to 2.7 times the rate of random selection
(\cref{tab:window}).

\cref{tab:recall_full} extends \cref{tab:recall} to all budgets with saved selections. Boundary scoring only
reorders the tiles inside the window, so the two values coincide at every $K$ outside the window. The coverage
view, and thus the prior, changes with the budget because $m=\lfloor K/2\rfloor$. Hits follow the definition of
\cref{sec:reliability}. In this table, the top of the refined ranking consists of the directly kept tiles in
prior order, followed by the window in the order of $q_t$ and the remaining tiles in prior order.

\paragraph{Why the top of the ranking needs a second stage.}
At $B=1536$, the first-ranked tile of the prior is the same as with the neutral question ``Describe this
image.'' on 90.6\% of the questions, and the same as with the swapped question on 91.6\%; the median
correlation between the prior with the real and the neutral question is 0.93. This tile falls into the top
quarter of the image on 63.4\% of the questions (62.3\%, 82.7\%, and 82.2\% at $B=1024$, $2048$, and $3072$), whereas
the top quarter holds the center of a target on only 4.7\% of the questions. The dependence on the question
therefore shows below the first rank, which is where the prior separates from the swapped-question prior
in \cref{fig:recall}. We have not established the cause of this pattern, and we did not change the method in
response to it.

\paragraph{Upper bound of the scorer.}
If $q_t$ is computed for all detail tiles, its first-ranked tile covers the target on 56.5\% of the questions,
and its top $k$ tiles on 84.8\%, against 71.2\% for the selection of \method at $B=1536$. This requires encoding
every tile and thus gives up the bound on compute, so it serves only as a reference. The neutral question
reaches a Recall@5 of 22.0\%, close to a center-first ranking (22.5\%), and the swapped question 49.2\%.

\paragraph{Target size.}
Measured in the coverage tiles at $B=1536$, 63.4\% of the questions have a target whose short side is below one
ViT patch (14\,px). Boundary scoring raises the hit rate from 59.5\% to 66.9\% on these questions, and from 77.1\%
to 78.6\% on the others.
